\origpage{41}{GRAPH MINORS. XIV}
\noindent \textit{chosen so that in addition $V(\mathscr{P}) \subseteq X$ and no member of $\mathscr{P}$ has any internal vertex in $\bar{\Omega}$.}

However, \del{we omit the proof, and} \fix{we} shall not use the result here.

\fix{We give the proof. In it we use the terms of [3], as in (8.2), and also the notation of [3], which is not related to the notation of (8.3). Thus if $(G \cup G_0, \Omega)$ is the composition of a society $(G_0, \Omega_0)$ with a neighbourhood $(G, \Omega, \Omega_0)$, then $G$ and $G_0$ have no common edge and $V(G \cap G_0) = \bar{\Omega}_0$. We shall prove the following.

\hspace*{12pt}(1) \textit{Let $(G, \Omega, \Omega_0)$ be a rural neighbourhood of breadth $\ge 3$, let $X \subseteq \bar{\Omega}$ be linked with $|X| \ge 4$, and let $(G \cup G_0, \Omega)$ be the composition of a society $(G_0, \Omega_0)$ with $(G, \Omega, \Omega_0)$. If $(G \cup G_0, \Omega)$ has a cross, then it has a cross $\mathscr{P}$ with $V(\mathscr{P}) \subseteq X$ such that no member of $\mathscr{P}$ has an internal vertex in $\bar{\Omega}$.}

This gives (8.3), applied to the neighbourhood $(G_1, \Omega_1, \Omega_2)$, the society $(G_2, \Omega_2)$, and the set $X'$ of initial vertices of the four given paths. Indeed $X' \subseteq X$ and $|X'| = 4$, since the paths are disjoint, and $X'$ is linked, by the four paths. Now (1) is the case $p = 2$ of [3, (10.4)], with breadth $\ge 3$ in place of $\ge p + 5$ and $|X| \ge 4$ in place of $|X| \ge 2p + 8$. The proof of [3, (10.4)] uses [3, (10.3)], whose proof uses [3, (5.2)], and only [3, (5.2)] needs $p + 4$ separating paths. We show first that for crosses two separating paths are enough, and then we follow the proofs of [3, (10.3)] and [3, (10.4)].

We use the following fact several times. Let $(H_1, H_2)$ be a separation of a graph, and let $R$ be a connected subgraph with a vertex in $V(H_1)$ and no vertex in $V(H_1 \cap H_2)$. Then $R \subseteq H_1$. For an edge of $R$ with an end in $V(H_1) - V(H_2)$ is not an edge of $H_2$, and so it is an edge of $H_1$, and its other end is in $V(H_1)$ and hence in $V(H_1) - V(H_2)$; since $R$ is connected, every edge and every vertex of $R$ is in $H_1$.

\hspace*{12pt}(2) \textit{Let $(G, \Omega)$ be a society with a cross. Let $\{Q_1, Q_2\}$ be a transaction in $(G, \Omega)$ such that $Q_1, Q_2$ are separating and $Q_i$ has ends $a_i, b_i$ $(i = 1, 2)$, where $(a_2, a_1, b_1, b_2)$ is clockwise. Let $(H_1, H_2)$ be a separation of $G$ with $V(H_1 \cap H_2) = V(Q_2)$, $Q_2 \subseteq H_1$, $\Omega(a_2, b_2) \subseteq V(H_1)$ and $\Omega(b_2, a_2) \subseteq V(H_2)$. Then either there is a cross $\{P, P'\}$ of $(G, \Omega)$ with $P, P' \subseteq H_1$, or there is a cross $\mathscr{P}$ of $(G, \Omega)$ with $\Omega(a_1, b_1) \cap V(\mathscr{P}) = \emptyset$.}

Since $(a_2, a_1, b_1, b_2)$ is clockwise, the vertices $a_2$ and $b_2$ are in $\Omega(b_1, a_1)$, and $\Omega(a_1, b_1) \subseteq \Omega(a_2, b_2)$. Suppose first that $Q_1, Q_2$ have no internal vertex in $\bar{\Omega}$. A cross is a crooked transaction of cardinality $2$, and so we may choose $\mathscr{P}$ and a fence for it such that the statement \textbf{S} of [3, Section 4] holds with $\mathscr{Q} = \{Q_1, Q_2\}$ and $p = 2$. Then $|\mathscr{P}| = 2$. For otherwise, by [3, (3.4)], some subset of $\mathscr{P}$ with two members is a crooked transaction; its union with $\bigcup \mathscr{Q}$ is included in $\bigcup \mathscr{P} \cup \bigcup \mathscr{Q}$, and so it is equal to it by (i) of \textbf{S}, and then (ii) of \textbf{S} is false. Hence $\mathscr{P}$ is a cross, and $|V(\mathscr{P})| = 4$. If $V(\mathscr{P}) \cap \Omega(a_1, b_1) = \emptyset$ the second outcome of (2) holds, and so we may assume that there is a vertex $x_0 \in V(\mathscr{P}) \cap \Omega(a_1, b_1)$.

\hspace*{12pt}(3) \textit{For $j = 1, 2$, if $P$ is an $\Omega$-path with both ends in $\Omega(a_j, b_j)$ and no vertex in $V(Q_j)$, then $P \subseteq H_1$.}

The ends of $P$ are in $\Omega(a_2, b_2) \subseteq V(H_1)$. We claim that $P$ has no vertex in $V(Q_2)$. For $j = 2$ this is a hypothesis. Let $j = 1$, and suppose that $P$ meets $V(Q_2)$. Let $R_1$ be a minimal subpath of $P$ between an end $u$ of $P$ and a vertex $w \in V(Q_2)$, and let $R_2$ be the subpath of $Q_2$ between $w$ and $a_2$. Then $R_1 \cup R_2$ is a path between $u$ and $a_2$, since $R_1$ meets $V(Q_2)$ only in $w$. It has no vertex in $V(Q_1)$, since $P$ has none and $Q_1, Q_2$ are disjoint. But $u \in \Omega(a_1, b_1)$ and $a_2 \in \Omega(b_1, a_1)$, and so $R_1 \cup R_2$ is an $\Omega$-path that crosses $Q_1$, which is impossible since $Q_1$ is separating. This proves our claim. Since $V(H_1 \cap H_2) = V(Q_2)$, the fact above shows that $P \subseteq H_1$. This proves (3).

Suppose that $Q_j \cap \bigcup \mathscr{P}$ is null for some $j \in \{1, 2\}$. Since $Q_j$ is separating and $x_0 \in \Omega(a_1, b_1) \subseteq \Omega(a_j, b_j)$, [3, (5.1)] shows that $\mathscr{P}$ has a subset which is a cross with all four ends in $\Omega(a_j, b_j)$. This subset is $\mathscr{P}$ itself, and by (3) both members of $\mathscr{P}$ are subgraphs of $H_1$, so the first outcome of (2) holds. We may therefore assume that $Q_1$ and $Q_2$ both meet $\bigcup \mathscr{P}$. Moreover $\mathscr{P}$ has no overpass, since removing a member of $\mathscr{P}$ leaves a transaction with one member, which does not have two peripheral members. So by [3, (4.3)], each of $Q_1, Q_2$ either has both ends in $V(\mathscr{P})$ or is half-covered.

\hspace*{12pt}(4) \textit{If $Q_j$ is half-covered, where $j \in \{1, 2\}$, then $\mathscr{P} = \{P, P'\}$, where $P$ and $Q_j$ have a common end $s$, $P \cap Q_j$ is a path between $s$ and some vertex $c$, $P'$ has no vertex in $V(Q_j)$, and the three vertices in $V(\mathscr{P}) - \{s\}$ are either all in $\Omega(a_j, b_j)$ or all in $\Omega(b_j, a_j)$.}

For let $P$ be the member of $\mathscr{P}$ that meets $Q_j$, as in [3, (4.3)(iii)], let $s$ be the common end of $P$, $Q_j$ and $P \cap Q_j$, and let $P'$ be the other member of $\mathscr{P}$. Then $P'$ has no vertex in $V(Q_j)$, since $Q_j$ meets only one member of $\mathscr{P}$. Since $Q_j$ is separating, $P'$ does not cross $Q_j$, and so both ends of $P'$ are in $\Omega(a_j, b_j)$ or both are in $\Omega(b_j, a_j)$. We assume the first; the second case is the same with the two sides exchanged. Let the ends of $P'$ be $u, w$, where $(a_j, u, w, b_j)$ is clockwise, and let $t$ be the end of $P$ other than $s$. Since $\mathscr{P}$ is crooked, $P'$ is not peripheral, and so $\Omega(u, w)$ contains $s$ or $t$. It does not contain $s$, since $s \in \{a_j, b_j\}$. Hence $t \in \Omega(u, w) \subseteq \Omega(a_j, b_j)$, and so $u, w, t$ are all in $\Omega(a_j, b_j)$. This proves (4).

Since $Q_2$ meets $\bigcup \mathscr{P}$, it has an end in $V(\mathscr{P})$, by (4) if $Q_2$ is half-covered. This end is in $\Omega(b_1, a_1)$, while $x_0 \in \Omega(a_1, b_1)$. So for $s \in \{a_1, b_1\}$ the set $V(\mathscr{P}) - \{s\}$ meets both $\Omega(a_1, b_1)$ and $\Omega(b_1, a_1)$, and hence $Q_1$ is not half-covered, by (4). Thus $a_1, b_1 \in V(\mathscr{P})$. If also $a_2, b_2 \in V(\mathscr{P})$, then $x_0, a_1, b_1, a_2, b_2$ are five distinct members of $V(\mathscr{P})$, which is impossible. Hence $Q_2$ is half-covered. Let $P, P', s, c$ be as in (4) with $j = 2$, and let $t$ be the end of $P$ other than $s$. Since $x_0 \in \Omega(a_2, b_2)$ and $x_0 \ne s$, the three vertices in $V(\mathscr{P}) - \{s\}$ are in $\Omega(a_2, b_2)$, by (4). In particular $P' \subseteq H_1$, by (3). Since $P \cap Q_2$ is the subpath of $P$ between $s$ and $c$, the subpath $R$ of $P$ between $c$ and $t$ meets $V(Q_2)$ only in $c$. Now $t \ne c$, since $t \in \Omega(a_2, b_2)$, while the only vertices of $Q_2$ in $\bar{\Omega}$ are $a_2$ and $b_2$. So $R \setminus c$ is connected, has no vertex in $V(Q_2) = V(H_1 \cap H_2)$, and contains $t \in V(H_1)$, and by the fact above $R \setminus c \subseteq H_1$. The edge of $R$ at $c$ is then in $H_1$, since its other end is in $V(H_1) - V(H_2)$, and $P \cap Q_2 \subseteq Q_2 \subseteq H_1$. Hence $P \subseteq H_1$, and the first outcome of (2) holds.

Now we drop the assumption that $Q_1, Q_2$ have no internal vertex in $\bar{\Omega}$. Let $G'$ be obtained from $G$ by adding, for each $v \in \bar{\Omega}$, a new vertex $v'$ and a new edge $vv'$, and let $\Omega'$ be the cyclic permutation with $\bar{\Omega}' = \{v' : v \in \bar{\Omega}\}$ that maps $v'$ to $w'$ whenever $\Omega$ maps $v$ to $w$. For an $\Omega$-path $P$ with ends $u, w$, let $P'$ be $P$ together with the edges $uu'$ and $ww'$. Then $P'$ is an $\Omega'$-path with no internal vertex in $\bar{\Omega}'$, and every $\Omega'$-path is $P'$ for exactly one $\Omega$-path $P$, since the new vertices have valency one. Two $\Omega$-paths are disjoint if and only if the corresponding $\Omega'$-paths are, and $\{u, w\}$ crosses $\{u_1, w_1\}$ in $\Omega$ if and only if $\{u', w'\}$ crosses $\{u_1', w_1'\}$ in $\Omega'$. Hence $P \mapsto P'$ and its inverse map crosses to crosses, transactions to transactions, and separating paths to separating paths. Let $H_1'$ be obtained from $H_1$ by adding the new edges at the vertices in $\Omega(a_2, b_2) \cup \{a_2, b_2\}$, with their ends, and let $H_2'$ be obtained from $H_2$ by adding the other new edges, with their ends, and the vertices $a_2', b_2'$. Then $(H_1', H_2')$ is a separation of $G'$ with $V(H_1' \cap H_2') = V(H_1 \cap H_2) \cup \{a_2', b_2'\} = V(Q_2')$ and $Q_2' \subseteq H_1'$, and $\Omega'(a_2', b_2') \subseteq V(H_1')$ and $\Omega'(b_2', a_2') \subseteq V(H_2')$. Since $(G', \Omega')$ has a cross, and $\{Q_1', Q_2'\}$ is a transaction of separating paths with no internal vertex in $\bar{\Omega}'$, the case already proved applies to $(G', \Omega')$, $\{Q_1', Q_2'\}$ and $(H_1', H_2')$. If it gives a cross of $(G', \Omega')$ whose members are subgraphs of $H_1'$, then the corresponding cross of $(G, \Omega)$ has members that are subgraphs of $H_1$, since $H_1'$ with the new vertices deleted is $H_1$. If it gives a cross $\mathscr{P}'$ with $\Omega'(a_1', b_1') \cap V(\mathscr{P}') = \emptyset$, the corresponding cross $\mathscr{P}$ of $(G, \Omega)$ has $\Omega(a_1, b_1) \cap V(\mathscr{P}) = \emptyset$. This proves (2).

Using (2) in place of [3, (5.2)], we follow the proof of [3, (10.3)] with two separating paths.

\hspace*{12pt}(5) \textit{Let $(G, \Omega, \Omega_0)$ be a rural neighbourhood of breadth $\ge 2$, let $X \subseteq \bar{\Omega}$ be weakly-linked with $|X| \ge 4$, and let $(G \cup G_0, \Omega)$ be the composition of a society $(G_0, \Omega_0)$ with $(G, \Omega, \Omega_0)$. If $(G \cup G_0, \Omega)$ has a cross, then it has a cross $\mathscr{P}$ with $V(\mathscr{P}) \subseteq X$.}

We show, by induction on $|Y - X|$, that for every $Y$ with $X \subseteq Y \subseteq \bar{\Omega}$, if $(G \cup G_0, \Omega)$ has a cross with all ends in $Y$ then it has one with all ends in $X$. The case $Y = \bar{\Omega}$ is (5). If $Y = X$ there is nothing to prove, and so we assume that there exists $x \in Y - X$. By the inductive hypothesis, applied to $Y - \{x\}$, it suffices to find a cross with all ends in $Y - \{x\}$. List the members of $X$ in clockwise order, starting just after $x$. Since $|X| \ge 4$, we may take $b_1, b_2$ to be the first two and $a_1, a_2$ to be the last two, and then $(a_2, a_1, x, b_1, b_2)$ is clockwise. The set $\{a_1, a_2, b_1, b_2\}$ is weakly-linked, since it is a subset of $X$, and so by [3, (10.2)] with $s = 2$ there is a transaction $\{Q_1, Q_2\}$ in $(G, \Omega)$ such that $Q_i$ has ends $a_i, b_i$ and meets every path of $G$ from $\Omega(a_i, b_i)$ to $\Omega(b_i, a_i) \cup \bar{\Omega}_0$ $(i = 1, 2)$.

Let $i \in \{1, 2\}$. Let $J$ be the union of the components of $G \setminus V(Q_i)$ that contain a vertex in $\Omega(a_i, b_i)$. Let $K_1$ consist of $J$, $Q_i$, and the edges of $G$ between $V(J)$ and $V(Q_i)$, and let $K_2$ consist of the vertices of $G$ not in $V(J)$ and the edges of $G$ not in $K_1$. An edge of $G$ with an end in $V(J)$ has its other end in $V(J) \cup V(Q_i)$, and so it is in $K_1$. Hence $(K_1, K_2)$ is a separation of $G$ with $V(K_1 \cap K_2) = V(Q_i)$, $Q_i \subseteq K_1$ and $\Omega(a_i, b_i) \subseteq V(K_1)$. Moreover $\Omega(b_i, a_i) \cup \bar{\Omega}_0 \subseteq V(K_2)$, since a vertex of $\Omega(b_i, a_i) \cup \bar{\Omega}_0$ in $V(J)$ would be joined to $\Omega(a_i, b_i)$ by a path of $G$ with no vertex in $V(Q_i)$. Since $V(K_1) \cap V(G_0) \subseteq \bar{\Omega}_0 \subseteq V(K_2)$, and $G_0$ has no edge in $G$, it follows that $(K_1, K_2 \cup G_0)$ is a separation of $G \cup G_0$ with $V(K_1 \cap (K_2 \cup G_0)) = V(Q_i)$. We call it the separation for $Q_i$.

Let $\Omega_Y$ be the cyclic permutation with $\bar{\Omega}_Y = Y$ that orders $Y$ as $\Omega$ does. The crosses of $(G \cup G_0, \Omega_Y)$ are the crosses of $(G \cup G_0, \Omega)$ with all ends in $Y$. Also $\{Q_1, Q_2\}$ is a transaction in $(G \cup G_0, \Omega_Y)$, since its ends are in $X \subseteq Y$, and $(a_2, a_1, b_1, b_2)$ is clockwise in $\Omega_Y$. For $i = 1, 2$ we have $\Omega_Y(a_i, b_i) \subseteq \Omega(a_i, b_i)$ and $\Omega_Y(b_i, a_i) \subseteq \Omega(b_i, a_i)$, and $Q_i$ is separating in $(G \cup G_0, \Omega_Y)$. For suppose that an $\Omega_Y$-path $R$ crosses $Q_i$, and let $(K_1, K_2 \cup G_0)$ be the separation for $Q_i$. Then $R$ has no vertex in $V(Q_i)$. It has one end in $\Omega(a_i, b_i) \subseteq V(K_1)$, and so $R \subseteq K_1$ by the fact above. But its other end is in $\Omega(b_i, a_i) - V(Q_i)$, which is disjoint from $V(K_1)$, a contradiction.

Now we apply (2) to $(G \cup G_0, \Omega_Y)$ and $\{Q_1, Q_2\}$, with $(H_1, H_2)$ the separation $(K_1, K_2 \cup G_0)$ for $Q_2$. Suppose first that (2) gives a cross $\{P, P'\}$ with $P, P' \subseteq K_1$. Then $\{P, P'\}$ is a cross of $(G, \Omega)$, since $K_1 \subseteq G$ and its ends are in $Y \subseteq \bar{\Omega}$. But $(G, \Omega)$ has no cross, since it is rural, as in the proof of [3, (10.3)], a contradiction. Hence (2) gives a cross $\mathscr{P}$ of $(G \cup G_0, \Omega_Y)$ with no end in $\Omega_Y(a_1, b_1)$. Its ends are in $Y$, and $x$ is not one of them since $x \in \Omega_Y(a_1, b_1)$. So $\mathscr{P}$ is a cross of $(G \cup G_0, \Omega)$ with all ends in $Y - \{x\}$. This proves (5).

Now we prove (1), following the proof of [3, (10.4)] with a $3$-nest in place of the $(p + 5)$-nest and (5) in place of [3, (10.3)]. Take a presentation $(\Sigma, \Gamma, \Delta, \Delta_0)$ of $(G, \Omega, \Omega_0)$ with a $3$-nest $(C_1', C_2', C_3')$, and let $\Delta_j'$ be the closed disc bounded by the drawing of $C_j'$ $(j = 1, 2, 3)$. The drawing of $C_2'$ is disjoint from that of $C_3'$, and so $\Delta_2'$ lies in the interior of $\Delta_3'$ and contains no point of $bd(\Delta)$. Hence no vertex of $C_1'$ or $C_2'$ is in $\bar{\Omega}$, and $\bar{\Omega}_0$, which is drawn in $\Delta_0 \subseteq \Delta_1'$, is disjoint from $\bar{\Omega}$. Let $n = |X|$. Since $X$ is linked, there are mutually disjoint $X \to \bar{\Omega}_0$ paths $P_1', \dots, P_n'$ of $G$, each with only its initial vertex in $\bar{\Omega}$. Let $G'$ be obtained from $G$ by deleting every edge of $G$ with an end in $\bar{\Omega}$ except the first edges of $P_1', \dots, P_n'$. Then $C_1', C_2'$ and $P_1', \dots, P_n'$ are subgraphs of $G'$, and $(\Sigma, \Gamma', \Delta, \Delta_0)$ is a presentation of $(G', \Omega, \Omega_0)$ with the $2$-nest $(C_1', C_2')$, where $\Gamma'$ is the drawing of $G'$ given by $\Gamma$. By [3, (10.1)] applied to $(G', \Omega, \Omega_0)$, there is a $2$-nest $(C_1, C_2)$ and there are disjoint $X \to \bar{\Omega}_0$ paths $P_1, \dots, P_n$ of $G'$, such that $P_i \cap C_j$ is a path for $1 \le i \le n$ and $j = 1, 2$. Each vertex $v \in \bar{\Omega}$ has valency at most $1$ in $G'$. For the edges of $G'$ at $v$ are first edges of the paths $P_k'$, the second vertex of each $P_k'$ is an internal vertex or is in $\bar{\Omega}_0$, and so it is not in $\bar{\Omega}$, and only one of these paths starts at $v$. Hence no vertex of $C_1$ or $C_2$ is in $\bar{\Omega}$, since circuits have no vertex of valency $1$. For the same reason no internal vertex of $P_i$ is in $\bar{\Omega}$, and so the only vertex of $P_i$ in $\bar{\Omega}$ is its initial vertex $x_i$, since its other end is in $\bar{\Omega}_0$.

Let $\Delta_1$ be the closed disc bounded by the drawing of $C_2$. As before, $\Delta_0$ lies in the interior of $\Delta_1$. No vertex in $\bar{\Omega}$ is drawn in $\Delta_1$. For such a vertex is drawn in $bd(\Delta)$, and it is not drawn in the drawing of $C_2$, since it is not a vertex of $C_2$, and so it would be in the interior of $\Delta_1$, which is included in the interior of $\Delta$. The drawing of $P_i$ runs from $x_i$, which is not in $\Delta_1$, to its end in $\bar{\Omega}_0$, which is in $\Delta_1$, and so it meets the drawing of $C_2$. Since an edge not in $C_2$ meets the drawing of $C_2$ only in its ends, $P_i$ has a vertex in $V(C_2)$; let $y_i$ be the first one, and put $Y = \{y_1, \dots, y_n\}$. Let $G_1$ be the subgraph of $G$ drawn in $\Delta_1$, and let $H_2$ be the subgraph of $G$ drawn in the closure of $\Delta - \Delta_1$, without the edges of $C_2$. An edge of $G$ not in $C_2$ is drawn in $\Delta_1$ or in the closure of $\Delta - \Delta_1$, and the vertices drawn in $bd(\Delta_1)$ are those of $C_2$. So $(G_1, H_2)$ is a separation of $G$ with $V(G_1 \cap H_2) = V(C_2)$. Since $V(G_0) \cap V(G) = \bar{\Omega}_0$, which is drawn in the interior of $\Delta_1$, it follows that $(G_1 \cup G_0, H_2)$ is a separation of $G \cup G_0$ with $V((G_1 \cup G_0) \cap H_2) = V(C_2)$, and $\bar{\Omega} \subseteq V(H_2)$. Let $\Omega_1$ be the cyclic permutation with $\bar{\Omega}_1 = V(C_2)$ given by the clockwise orientation of $bd(\Delta_1)$. Then $(G_1, \Omega_1, \Omega_0)$ and $(H_2, \Omega, \Omega_1)$ are rural neighbourhoods, the first with the $2$-nest $(C_1, C_2)$, and $(G \cup G_0, \Omega)$ is the composition of $(G_1 \cup G_0, \Omega_1)$ with $(H_2, \Omega, \Omega_1)$. By [3, (9.1)] and [3, (3.4)], $(G_1 \cup G_0, \Omega_1)$ has a cross.

The set $Y$ is weakly-linked in $(G_1, \Omega_1, \Omega_0)$, and $|Y| = n \ge 4$. For let $1 \le i \le n$. The path $P_i \cap C_2$ contains $y_i$, and it is a subpath of $P_i$, and so $y_i$ is one of its ends, since $y_i$ is the first vertex of $P_i$ in $V(C_2)$. Let $z_i$ be its other end, and let $T_i$ be the subpath of $P_i$ from $z_i$ to the end $t_i$ of $P_i$ in $\bar{\Omega}_0$. Then $T_i$ meets $V(C_2)$ only in $z_i$, and $t_i \ne z_i$, since $t_i$ is drawn in the interior of $\Delta_1$. So $T_i \setminus z_i$ is a connected subgraph of $G$ with no vertex in $V(C_2)$ and with $t_i \in V(G_1)$, and it is a subgraph of $G_1$ by the fact above, applied to $(G_1, H_2)$. The edge of $T_i$ at $z_i$ is in $G_1$ as well, since its other end is in $V(G_1) - V(C_2)$. Hence the subpath of $P_i$ from $y_i$ to $t_i$, which is the union of $P_i \cap C_2$ and $T_i$, is a path of $G_1$. These paths are disjoint, since $P_1, \dots, P_n$ are, and so $Y$ is weakly-linked.

By (5), applied to $(G_1, \Omega_1, \Omega_0)$, $Y$ and $(G_1 \cup G_0, \Omega_1)$, there is a cross $\{P, P'\}$ of $(G_1 \cup G_0, \Omega_1)$ with all ends in $Y$. Let $y_{i_1}, y_{i_2}$ be the ends of $P$ and $y_{i_3}, y_{i_4}$ those of $P'$, and for $m = 1, \dots, 4$ let $R_m$ be the subpath of $P_{i_m}$ between $x_{i_m}$ and $y_{i_m}$. Then $R_m$ meets $V(C_2)$ only in $y_{i_m}$, by the choice of $y_{i_m}$, and so its drawing meets the drawing of $C_2$ only in $y_{i_m}$. Since $x_{i_m}$ is not drawn in $\Delta_1$, the rest of $R_m$ is drawn outside $\Delta_1$. So $R_m$ is a path of $H_2$ from $\bar{\Omega}$ to $\bar{\Omega}_1$, and it meets $G_1 \cup G_0$ only in $y_{i_m}$, since $V(G_0) \cap V(G) = \bar{\Omega}_0 \subseteq V(G_1)$. Let $P^+ = P \cup R_1 \cup R_2$ and $P'^+ = P' \cup R_3 \cup R_4$. These are disjoint paths of $G \cup G_0$, with ends $x_{i_1}, x_{i_2}$ and $x_{i_3}, x_{i_4}$ in $X$, since $P, P'$ are disjoint and $R_1, \dots, R_4$ are disjoint. No internal vertex of $P^+$ or $P'^+$ is in $\bar{\Omega}$, since no vertex of $G_1 \cup G_0$ is in $\bar{\Omega}$ and $x_{i_m}$ is the only vertex of $R_m$ in $\bar{\Omega}$. The paths $R_1, \dots, R_4$ are mutually disjoint, and they are drawn in the closed annulus which is the closure of $\Delta - \Delta_1$, each from one of its boundary circles to the other. So the ends $x_{i_1}, \dots, x_{i_4}$ occur around $bd(\Delta)$ in the same cyclic order as the ends $y_{i_1}, \dots, y_{i_4}$ occur around $bd(\Delta_1)$. Since $\{y_{i_1}, y_{i_2}\}$ crosses $\{y_{i_3}, y_{i_4}\}$ in $\Omega_1$, it follows that $\{x_{i_1}, x_{i_2}\}$ crosses $\{x_{i_3}, x_{i_4}\}$ in $\Omega$. Thus $\{P^+, P'^+\}$ is a cross as required in (1). This proves (1), and hence (8.3).}\fixnote{fix:8.3-proof-added}{The paper omits the proof of (8.3). The proof above is the one used in the formalization.}

The main result of this section is the following.

\stmt
(8.4) \textit{Let $\Sigma$ be a closed cylinder with ends $F_1, F_2$, and let $H$ be a rigid drawing in $\Sigma$. Let $C_1, C_2$ be disjoint circuits of $H$, both winding around $\Sigma$, and such that $F_2$ and $U(C_2)$ both lie in the same component of $\Sigma - U(C_1)$. Let $\pi \ge 0$ with $\pi \ne 1, 3$ and let the regional distance between $U(C_1)$ and $U(C_2)$ be at least $\pi + 4$. Let $X \subseteq V(C_1)$, and let there be $2\pi + 8$ mutually disjoint paths of $H$ from $X$ to $V(C_2)$, each with no vertex in $C_1$ except its initial vertex. Let $H$ be a subgraph of $G_1$, let $\mathscr{S}$ be a 3-segregation of $G_1$, and let $\alpha$ be an arrangement of $\mathscr{S}$ in $\Sigma$ compatible with $H$. Let $\Omega_i = \alpha^{-1}(F_i) \, (i = 1, 2)$, and let $\Omega_0$ be a natural order of $V(C_1)$ from $U(C_1)$. Now let $(G, \Omega_1)$ be a society which is the composition of a society $(G_2, \Omega_2)$ and the neighbourhood $(G_1, \Omega_1, \Omega_2)$, and let $(G, \Omega_1)$ have a crooked transaction of cardinality $\ge \pi$. Then $(G, \Omega_0)$ has a crooked transaction $\mathscr{P}$ of cardinality $\ge \pi$ such that for each $P \in \mathscr{P}$ with ends $u, v$ say, $u, v \in X$ and $U(P \cap H) - \{u, v\}$ is included in the component of $\Sigma - U(C_1)$ which includes $U(C_2)$.}

\textit{Proof.} We proceed by induction on $|V(G_1)| + |E(G_1)|$. We claim that we may assume that

\hspace*{12pt}(1) \textit{For each $(A, \Omega) \in \mathscr{S}$, there is no proper subgraph $A' \subseteq A$ such that $\bar{\Omega} \subseteq V(A')$, $H \cap A \subseteq A'$, $E(A') \ne \emptyset$ and for all $u, v \in \bar{\Omega}$, if $u$ and $v$ are joined by a path of $A$ with no internal vertex in $\bar{\Omega}$ then they are joined by such a path in $A'$.}

\textit{Subproof.} Suppose that such $A'$ exists, and let $\mathscr{S}'$ be obtained from $\mathscr{S}$ by replacing $(A, \Omega)$ by $(A', \Omega)$ (since $E(A') \ne \emptyset$, $(A', \Omega)$ is not already a member of $\mathscr{S}$ and so $\mathscr{S}'$ is still a set rather than a ``multiset''). Let $G'_1$ be the graph such that $\mathscr{S}'$ is a 3-segregation of $G'_1$, and let $G' = G'_1 \cup G_2$. We claim that $(G', \Omega_1)$ has a crooked transaction of order $\ge \pi$. To see this, let $\mathscr{P}$ be a crooked transaction in $(G, \Omega_1)$ of order $\ge \pi$; then we may assume that $\mathscr{P}$ is not a crooked transaction in $(G', \Omega_1)$, and so there exists $P \in \mathscr{P}$ with $P \cap A \not\subseteq A'$. Since $V(\mathscr{P}) \subseteq \bar{\Omega}_1 \subseteq V(\mathscr{S})$, at least two vertices of $P$ are in $\bar{\Omega}$; and hence $P$ is unique, because $|\bar{\Omega}| \le 3$. There is a unique subpath $Q$ of $P$ with ends in $\bar{\Omega}$, with $Q \subseteq A$ and with $Q \not\subseteq A'$; and it is easy to check that, from the assumptions about $A'$ in (1), there is a path $Q'$ of $A'$ with the same ends as $Q$ and with $\bar{\Omega} \cap V(Q') \subseteq \bar{\Omega} \cap V(Q)$. If we replace the subpath $Q$ of $P$ by $Q'$, we obtain a crooked transaction in $(G', \Omega_1)$ of order $\ge \pi$, as claimed. \fix{We give the details, replacing $P$ as a whole. Let $B$ be the subgraph of $G$ such that $(A, B)$ is a separation of $G$ with $V(A \cap B) = \bar{\Omega}$. Then $V(\mathscr{P}) \subseteq V(B)$, since $V(A) \cap V(\mathscr{S}) = \bar{\Omega}$. We claim that for every path $R$ of $G$ with both ends in $V(B)$, there is a path $R'$ of $A' \cup B$ with the same ends such that $V(R') \cap V(B) \subseteq V(R) \cap V(B)$. We prove this by induction on $|V(R)|$. If $R \subseteq B$ we take $R' = R$. If some internal vertex $w$ of $R$ is in $\bar{\Omega}$, we apply induction to the two subpaths of $R$ between its ends and $w$, and take for $R'$ a path between the ends of $R$ in the union of the two paths obtained. Otherwise $R \not\subseteq B$ and no internal vertex of $R$ is in $V(A \cap B)$, and so $R \subseteq A$ and both ends of $R$ are in $\bar{\Omega}$. By the assumption about $A'$ in (1), there is a path $R'$ of $A'$ with the same ends as $R$ and no internal vertex in $\bar{\Omega}$, and then $V(R') \cap V(B) \subseteq V(R) \cap V(B)$, as required. This proves our claim. Now every member of $\mathscr{P}$ that is not a subgraph of $B$ has at least two vertices in $\bar{\Omega}$, and so $P$ is the only such member, since $|\bar{\Omega}| \le 3$. Let $P'$ be a path of $A' \cup B \subseteq G'$ obtained from $P$ by our claim. Then $P'$ is disjoint from the other members of $\mathscr{P}$, since they are subgraphs of $B$ and $V(P') \cap V(B) \subseteq V(P)$. Hence replacing $P$ by $P'$ gives a crooked transaction in $(G', \Omega_1)$ with the same ends.}\fixnote{fix:8.4-claim1-step-added}{The paper says only that it is easy to check that the path $Q'$ exists. The argument above is the one used in the formalization.} But then from the inductive hypothesis, since $H \subseteq G'_1$, the result holds for $G'$ and hence for $G$, as required. This proves that we may assume (1).
